% IJCIP abstract: one paragraph, at most 250 words, no citations, no equations,
% no undefined abbreviations. Verified by experiments/check_submission.py.
\newcommand{\capmabstract}{%
Vehicle cybersecurity research and regulation concentrate on the product:
standards for road-vehicle cybersecurity engineering and type approval govern
the car, not the factory that builds it. Yet the plant is the asset whose
failure halts output and propagates through a continental supply chain. The
2025 attack on one manufacturer stopped global production for five weeks and
was assessed as costing a national economy about 1.9 billion pounds. This paper
treats the automotive plant as critical infrastructure and asks by which paths
an attacker reaches the assembly line. We contribute a curated
corpus of twenty publicly reported incidents in automotive manufacturing
between 2017 and 2025, each reconstructed as an asset-level path with an
explicit confidence label, and a cross-plane attack-path model spanning
corporate information technology, the identity plane, cloud and
software-as-a-service, the industrial demilitarised zone, manufacturing
execution, supervisory control, controllers and the supplier interface. The
model is instantiated as a reference plant of twelve zones, forty-eight assets
and ninety-one attack steps drawn from forty-three techniques of the MITRE
adversarial tactics and techniques knowledge bases for enterprise and
industrial systems, and analysed by exact most-likely-path ranking and Monte
Carlo simulation of attacker campaigns. Paths that never descend below the enterprise
level are 8.5 percent of the graph yet carry 68.7 percent of the aggregate
likelihood. Recovery capability dominates expected loss, controls that appear
worthless alone prove the least removable, and partial hardening displaces risk
between adversary objectives rather than removing it.%
}

\newcommand{\capmkeywords}{%
critical infrastructure protection \sep automotive manufacturing \sep
attack graphs \sep operational technology \sep cascading failure \sep
supply-chain security \sep risk quantification%
}
